
% ---------- Exercise 13 ----------
\item {\bfseries\boldmath
    State, with justification, whether the following
    statements are true or false.
}

Let $n\in\mathbb{N}^+$ and
$A,B,C\in\mathcal{M}_n(\mathbb{K})$.

\begin{enumerate}[label=\alph*)]

% ---------- a) ----------
\item {\bfseries\boldmath
    For any square matrices $A$ and $B$
    of the same order, $AB=BA$.
}

\begin{answer}
    \textbf{Universal statement:}

    $$
    \forall A,B\in\mathcal{M}_n(\mathbb{K}):
    AB=BA.
    $$

    \textbf{False. Proof by counterexample.}

    Consider

    $$
    A=\begin{bmatrix}
        1 & 0 \\
        0 & 0
    \end{bmatrix},
    \qquad
    B=\begin{bmatrix}
        0 & 1 \\
        0 & 0
    \end{bmatrix}.
    $$

    Then,

    $$
    AB=\begin{bmatrix}
        0 & 1 \\
        0 & 0
    \end{bmatrix}
    \neq
    \begin{bmatrix}
        0 & 0 \\
        0 & 0
    \end{bmatrix}=BA.
    $$

    Therefore, the statement is false.

    \par\bigskip
    \noindent\textbf{Sufficient condition:}

    If $A$ and $B$ are diagonal matrices,
    then, for all $i,j\in\{1,\ldots,n\}$,

    $$
    \begin{aligned}
        (AB)_{ii}
            &=a_{ii}b_{ii}
            =b_{ii}a_{ii}
            =(BA)_{ii},\\[0.5em]
        (AB)_{ij}
            &=0=(BA)_{ij},
            \quad i\neq j.
    \end{aligned}
    $$

    Hence, all corresponding entries
    are equal, so

    $$
    \boxed{AB=BA}.
    $$
\end{answer}


% ---------- b) ----------
\item {\bfseries\boldmath
    For any square matrices $A$ and $B$
    of the same order, if $AB=0$,
    then $A=0$ or $B=0$.
}

\begin{answer}
    \textbf{Universal statement:}

    $$
    \forall A,B\in\mathcal{M}_n(\mathbb{K}):
    AB=0_n\implies
    (A=0_n\lor B=0_n).
    $$

    \textbf{False. Proof by counterexample.}

    Consider

    $$
    A=B=\begin{bmatrix}
        0 & 1 \\
        0 & 0
    \end{bmatrix}
    \neq0_2.
    $$

    However,

    $$
    AB=\begin{bmatrix}
        0 & 1 \\
        0 & 0
    \end{bmatrix}
    \begin{bmatrix}
        0 & 1 \\
        0 & 0
    \end{bmatrix} =\begin{bmatrix}
        0 & 0 \\
        0 & 0
    \end{bmatrix}=0_2
    $$

    Thus,

    $$
    AB=0_2
    \quad\land\quad
    A\neq0_2
    \quad\land\quad
    B\neq0_2.
    $$

    Therefore, the statement is false.

    \par\bigskip
    \noindent\textbf{Sufficient condition:}

    If either $A$ or $B$ is invertible,
    the implication holds.

    \par\smallskip
    \noindent\textit{Case 1: $A$ is invertible.}

    Multiplying by $A^{-1}$ on the left:

    $$
    \begin{aligned}
        AB=0_n
        &\implies A^{-1}(AB)=A^{-1}0_n \\
        &\implies (A^{-1}A)B=0_n \\
        &\implies I_nB=0_n \\
        &\implies \boxed{B=0_n}.
    \end{aligned}
    $$

    \par\smallskip
    \noindent\textit{Case 2: $B$ is invertible.}

    Multiplying by $B^{-1}$ on the right:

    $$
    \begin{aligned}
        AB=0_n
        &\implies (AB)B^{-1}=0_nB^{-1} \\
        &\implies A(BB^{-1})=0_n \\
        &\implies AI_n=0_n \\
        &\implies \boxed{A=0_n}.
    \end{aligned}
    $$

    Hence, if at least one factor is
    invertible, the other must be zero
    whenever $AB=0_n$.
\end{answer}


% ---------- c) ----------
\item {\bfseries\boldmath
    For any square matrices $A$, $B$,
    and $C$ of the same order, if $AB=AC$,
    then $B=C$.
}

\begin{answer}
    \textbf{Universal statement:}

    $$
    \forall A,B,C\in\mathcal{M}_n(\mathbb{K}):
    AB=AC\implies B=C.
    $$

    \textbf{False. Proof by counterexample.}

    Consider

    $$
    A=\begin{bmatrix}
        1 & 0 \\
        0 & 0
    \end{bmatrix},
    \qquad
    B=\begin{bmatrix}
        1 & 0 \\
        0 & 1
    \end{bmatrix},
    \qquad
    C=\begin{bmatrix}
        1 & 0 \\
        0 & 0
    \end{bmatrix}.
    $$

    Then,

    $$
    AB=\begin{bmatrix}
        1 & 0 \\
        0 & 0
    \end{bmatrix}
    =AC,
    $$

    while $B\neq C$.

    Therefore, the statement is false.

    \par\bigskip
    \noindent\textbf{Sufficient condition:}

    If $A$ is invertible, there exists
    $A^{-1}\in\mathcal{M}_n(\mathbb{K})$
    such that

    $$
    A^{-1}A=AA^{-1}=I_n.
    $$

    Multiplying both sides of $AB=AC$
    by $A^{-1}$ on the left:

    $$
    \begin{aligned}
        AB=AC
        &\implies A^{-1}(AB)=A^{-1}(AC) \\
        &\implies (A^{-1}A)B=(A^{-1}A)C \\
        &\implies I_nB=I_nC \\
        &\implies \boxed{B=C}.
    \end{aligned}
    $$

    Hence,

    $$
    \boxed{
        A\text{ invertible}
        \land AB=AC
        \implies B=C.
    }
    $$
\end{answer}


% ---------- d) ----------
\item {\bfseries\boldmath
    If, for a square matrix $A$, there exists
    $k\in\mathbb{N}$ such that $A^k=0$,
    then $A=0$.
}

\begin{answer}
    \textbf{Universal statement:}

    $$
    \forall A\in\mathcal{M}_n(\mathbb{K}):
    \left(
        \exists k\in\mathbb{N}^+:
        A^k=0_n
    \right)
    \implies A=0_n.
    $$

    \textbf{False. Proof by counterexample.}

    Consider

    $$
    A=\begin{bmatrix}
        0 & 1 \\
        0 & 0
    \end{bmatrix}
    \neq0_2.
    $$

    For $k=2$,

    $$
    A^2
    =\begin{bmatrix}
        0 & 1 \\
        0 & 0
    \end{bmatrix}
    \begin{bmatrix}
        0 & 1 \\
        0 & 0
    \end{bmatrix} =0_2
    $$

    Thus,

    $$
    \exists k=2\in\mathbb{N}^+:
    A^k=0_2,
    \quad A\neq0_2.
    $$

    Therefore, the statement is false.

    \par\bigskip
    \noindent\textbf{Sufficient condition:}

    Suppose $A$ is idempotent, i.e.,

    $$
    A^2=A.
    $$

    By mathematical induction,

    $$
    \forall p\in\mathbb{N}^+:
    A^p=A.
    $$

    Indeed, the base case $A^1=A$ holds.
    Assuming $A^p=A$ for some $p\geq1$,
    multiplying by $A$ gives

    $$
    \begin{aligned}
        A^{p+1}
        &=A^pA \\
        &=AA \\
        &=A^2 \\
        &=A.
    \end{aligned}
    $$

    Now, if

    $$
    \exists k\in\mathbb{N}^+:
    A^k=0_n,
    $$

    then, since $A^k=A$,

    $$
    A=A^k=0_n.
    $$

    Therefore,

    $$
    \boxed{A=0_n}.
    $$

    A matrix that is both idempotent and nilpotent must be zero.
\end{answer}


% ---------- e) ----------
\item {\bfseries\boldmath
    For any square matrices $A$ and $B$
    of the same order, $(AB)^2=A^2B^2$.
}

\begin{answer}
    \textbf{Universal statement:}

    $$
    \forall A,B\in\mathcal{M}_n(\mathbb{K}):
    (AB)^2=A^2B^2.
    $$

    \textbf{False. Proof by counterexample.}

    Consider

    $$
    A=\begin{bmatrix}
        0 & 1 \\
        0 & 0
    \end{bmatrix},
    \qquad
    B=\begin{bmatrix}
        0 & 0 \\
        1 & 0
    \end{bmatrix}.
    $$

    Computing $AB$,

    $$
    AB=\begin{bmatrix}
        1 & 0 \\
        0 & 0
    \end{bmatrix}.
    $$

    Therefore,

    $$
    (AB)^2=\begin{bmatrix}
        1 & 0 \\
        0 & 0
    \end{bmatrix}.
    $$

    However,

    $$
    A^2=0_2,
    \qquad B^2=0_2,
    $$

    so

    $$
    A^2B^2=0_2.
    $$

    Hence,

    $$
    (AB)^2\neq A^2B^2,
    $$

    and the statement is false.

    \par\bigskip
    \noindent\textbf{Sufficient condition:}

    Suppose $AB=BA$.

    Then,

    $$
    \begin{aligned}
        (AB)^2
        &=(AB)(AB) \\
        &=A(BA)B \\
        &=A(AB)B \\
        &=(AA)(BB) \\
        &=\boxed{A^2B^2}.
    \end{aligned}
    $$

    Therefore,

    $$
    \boxed{
        AB=BA\implies(AB)^2=A^2B^2.
    }
    $$
\end{answer}

\end{enumerate}
